\section{Discussion}
\label{sec:discussion}

\subsection{How to Select a Design}
\label{sec:select}

The limits of \cref{sec:limits} tell when each design of \cref{tab:designs} is the correct
selection.

\begin{itemize}
  \item Use QO if each order can wait until a quorum signs it. Then the hostile tenure is
  zero, and no other design is safer.
  \item Use SL if no officer can be hostile, and the attacker cannot get the key of the
  server. Then one server is sufficient.
  \item Use NE if each order is easy to cancel and the system has a record of all orders.
  Then a long hostile tenure causes small damage.
  \item Use \sys{} if the commander must continue through outages, an officer can be
  hostile, and orders can cause damage that is not easy to repair.
\end{itemize}

A designer who selects \sys{} makes three decisions. The first decision is the horizon
$H$. It is the longest outage that a loyal commander must continue through. It is also the
time that a hostile commander keeps a device, and no protocol can change this
(\cref{thm:limit}). Thus the designer compares two risks: a time $H$ with a hostile
commander, and an outage longer than $H$ with only local control.

The second decision is the fallback period $F$. With $F = 0$, the system has unity of
command, and it has no commander between the end of the last renewal and a new grant. With
$F > 0$, the deputy can give orders in that period, and a hostile deputy gets the same
period at each device.

The third decision is the quorum $q$. The value $q = \lceil (n + f + 1)/2 \rceil$ is the
minimum. A larger $q$, to a maximum of $n - f$, makes the two thresholds of
\cref{cor:thresholds} smaller. Then a smaller group of loyal officers can remove a hostile
commander. But a smaller group of loyal officers, together with the hostile officers, can
also stop the renewals of a loyal commander.

\subsection{Limits of the Model}
\label{sec:model-limits}

\paragraph{Clocks}
The tenure limit uses the clock of the device. An attacker who sets that clock to a previous time extends the tenure. Substations frequently get time from satellite signals, and an
attacker can forge such signals~\cite{zhang2013tsa,gao2023quick}. Methods to find such an
attack are known~\cite{laverty2022gnss}. A device must thus keep a local clock, and
$\eps$ must include the error of that clock during the longest outage. \pending{Not done at
this time: a model of clock attacks.}

\paragraph{More than $f$ hostile officers}
If more than $f$ officers are hostile, it is possible that each officer who is in two
quorums is hostile. Then \cref{lem:grant} and \cref{lem:separation} are not correct, and two commanders
are possible. With $q$ hostile officers, no property stays.

\paragraph{A constant set of officers}
The set of officers and their keys do not change in the model. A change of that set is a
different problem~\cite{kuznetsov2022reconfig}. We did not examine it.

\paragraph{Hostile devices}
Each device is correct in the model. A device with changed software can do each order and
can ignore each rule. The report on the 2015 attack records an attack on the firmware of
network devices at substations~\cite[p.~9]{eisac2016ukraine}.

\paragraph{The decision to stop consent}
The model does not tell how an officer knows that a commander is hostile. A second
channel, for example a telephone, helps the officers to make that decision. An attacker can
also attack that channel. The telephone calls of the 2015 attack are an
example~\cite[p.~2]{eisac2016ukraine}.

\paragraph{Damage in the tenure}
\sys{} puts a limit on the time of a hostile commander, and not on the damage. In that
time, the hostile commander can give each order. Checks of the contents of
orders~\cite{mashima2018onthefly,hong2019resilient} and a limit on the number of orders
decrease the damage. \Cref{thm:limit} shows that such checks cannot make the tenure
shorter.

\paragraph{A faster stop when messages come through}
A device that receives a valid grant of a higher epoch can reject the orders of the
previous epoch immediately. This rule is safe, because a quorum signs each valid grant. It
makes the tenure shorter when the attacker does not stop the grant. \Cref{thm:limit} shows
that a design cannot use this rule as a replacement for expiry.

\paragraph{An attacker who only stops messages}
An attacker who controls the links can always stop orders and renewals. Expiry thus gives
that attacker no new method to stop command. A takeover on a timeout is different: it
gives the attacker a method to move command.

\subsection{Known Risks of Errors in This Draft}

Five risks are known. First, no second person and no model checker examined the proofs.
Second, the data about related work come from abstracts, and \cref{tab:position} can
be incorrect for a source. Third, the entries of \cref{tab:designs} for the four known
designs come from definitions and not from tests. Fourth, no power engineer examined the
grid example and rule D3. Fifth, the numbers of \cref{sec:baseline} come from a model of
incident response with no attacker.
